\documentclass[11pt]{scrartcl}
\usepackage[mdthm,colorsec,noauthor,noasy]{evan}
\DeclareMathOperator{\nul}{null}
\DeclareMathOperator{\range}{range}
\DeclareMathOperator{\spn}{span}
\newcommand{\LL}{\mathcal L}
\newcommand{\PP}{\mathcal P}
\newcommand{\MM}{\mathcal M}
\title{Math 110 Hard Mock Midterm}
\subtitle{50 minutes \quad 4 questions \quad 40 points}
\author{}
\date{}
\ihead{}
\begin{document}
\maketitle

\noindent Set a 50-minute timer. Give concise reasoning with each answer. You may quote results from the notes or text if you state them. $\PP_n(\RR)$ is the space of real polynomials of degree at most $n$. All bases are ordered.

\medskip\noindent\textbf{1. Is it direct? (8 points)} In $\PP_3(\RR)$ let
\[ U=\{p:p(1)=p(-1)=0\},\qquad W=\spn\big(1+x^3,\ x+x^2\big). \]
\begin{enumerate}[a)]
  \ii (4) Determine whether $\PP_3(\RR)=U\oplus W$. Prove your answer.
  \ii (2) Find $\dim(U+W)$ and a basis of $U\cap W$.
  \ii (2) Find a $2$-dimensional subspace $W'$ containing $1+x^3$ with $\PP_3(\RR)=U\oplus W'$, and prove it works.
\end{enumerate}

\medskip\noindent\textbf{2. When null space and range split the space (10 points)} Let $V$ be finite-dimensional and $T\in\LL(V)$.
\begin{enumerate}[a)]
  \ii (2) Prove that $\nul T\subseteq\nul T^2$.
  \ii (3) Let $R$ be $T$ restricted to $\nul T^2$. Prove that $\range R=\nul T\cap\range T$ and $\nul R=\nul T$. Deduce
  \[ \dim\nul T^2=\dim\nul T+\dim(\nul T\cap\range T). \]
  \ii (3) Prove that the following are equivalent: (i) $\nul T^2=\nul T$; (ii) $\nul T\cap\range T=\{0\}$; (iii) $V=\nul T\oplus\range T$.
  \ii (2) Let $D\colon\PP_3(\RR)\to\PP_3(\RR)$ be differentiation. Use (c) to decide whether $\PP_3(\RR)=\nul D\oplus\range D$.
\end{enumerate}

\medskip\noindent\textbf{3. A matrix in a nonstandard basis (10 points)} Define $T\colon\PP_2(\RR)\to\PP_2(\RR)$ by
\[ (Tp)(x)=p(x+1)-x\,p'(x). \]
Let $B=\big(1,\ x,\ x(x-1)\big)$.
\begin{enumerate}[a)]
  \ii (4) Find $\MM(T)$ with respect to $B$ on both sides.
  \ii (3) Find a basis of $\nul T$ and a basis of $\range T$. Is $T$ an isomorphism?
  \ii (3) Find all $p\in\PP_2(\RR)$ with $Tp=1+2x-x^2$.
\end{enumerate}

\medskip\noindent\textbf{4. Operators commuting with a projection (12 points)} Let $V$ be finite-dimensional with $\dim V=n$, and let $P\in\LL(V)$ satisfy $P^2=P$. Let $k=\dim\range P$.
\begin{enumerate}[a)]
  \ii (3) Prove that $V=\nul P\oplus\range P$.
  \ii (4) Let $T\in\LL(V)$. Prove that $TP=PT$ if and only if $T(\nul P)\subseteq\nul P$ and $T(\range P)\subseteq\range P$.
  \ii (3) Prove that $\mathcal C=\{T\in\LL(V):TP=PT\}$ is a subspace of $\LL(V)$ and that $\dim\mathcal C=k^2+(n-k)^2$.
  \ii (2) Let $P(x,y,z)=(x,y,0)$ on $\RR^3$. Find $\dim\mathcal C$, and give an operator on $\RR^3$ that is not in $\mathcal C$.
\end{enumerate}

\newpage
\section*{Solutions}

\subsection*{Question 1}
\textbf{(a)} Not direct. Every $p=a(1+x^3)+b(x+x^2)\in W$ has $p(-1)=0$ automatically, and $p(1)=2a+2b$. Taking $a=1$, $b=-1$ gives
\[ p=1-x-x^2+x^3=(1-x)^2(1+x)\ne0, \]
which lies in $W$ and vanishes at $\pm1$, so it lies in $U$. Hence $U\cap W\ne\{0\}$ and the sum is not direct.
(Trap: $\dim U=2$ and $\dim W=2$ add to $4$, but the dimension count proves nothing until the intersection is checked.)

\textbf{(b)} $U$ is the null space of $p\mapsto(p(1),p(-1))$, which is surjective ($\tfrac12(1+x)\mapsto(1,0)$, $\tfrac12(1-x)\mapsto(0,1)$), so $\dim U=2$. $1+x^3$ and $x+x^2$ are independent, so $\dim W=2$. The computation in (a) shows $U\cap W=\spn\big(1-x-x^2+x^3\big)$, dimension $1$. So $\dim(U+W)=2+2-1=3$.

\textbf{(c)} $W'=\spn(1+x^3,\,1)$. If $p=a(1+x^3)+b\in U$, then $p(-1)=b=0$ and $p(1)=2a+b=0$, so $a=b=0$. So $U\cap W'=\{0\}$, and $2+2=4=\dim\PP_3$, so $\PP_3=U\oplus W'$.

\subsection*{Question 2}
\textbf{(a)} If $Tv=0$ then $T^2v=T(0)=0$.

\textbf{(b)} Range: if $v\in\nul T^2$ then $T(Tv)=0$, so $Tv\in\nul T$, and $Tv\in\range T$. Conversely if $y\in\nul T\cap\range T$, write $y=Tv$; then $T^2v=Ty=0$, so $v\in\nul T^2$ and $Rv=y$.
Null space: $\nul R=\nul T\cap\nul T^2=\nul T$ by (a).
Rank--nullity for $R$ on $\nul T^2$ gives the formula.

\textbf{(c)} (i)$\Leftrightarrow$(ii): by (a), $\nul T\subseteq\nul T^2$, so they are equal exactly when the dimensions agree, which by (b) is exactly when $\dim(\nul T\cap\range T)=0$.
(ii)$\Rightarrow$(iii): the sum is direct, and $\dim\nul T+\dim\range T=\dim V$ by rank--nullity, so it is all of $V$.
(iii)$\Rightarrow$(ii): a direct sum of two subspaces has zero intersection.

\textbf{(d)} $\nul D=\spn(1)$ and $\nul D^2=\PP_1\ne\nul D$. By (c), $\PP_3\ne\nul D\oplus\range D$. (Indeed $1\in\nul D\cap\range D$.)

\subsection*{Question 3}
\textbf{(a)} $T(1)=1$. $T(x)=(x+1)-x=1$. $T(x^2)=(x+1)^2-2x^2=-x^2+2x+1$, so $T\big(x(x-1)\big)=T(x^2)-T(x)=-x^2+2x=-x(x-1)+x$. In the basis $B$:
\[ \MM(T)=\begin{pmatrix}1&1&0\\0&0&1\\0&0&-1\end{pmatrix}. \]
(Check one column a second way: $x(x-1)$ has $p(x+1)=x^2+x$ and $xp'=2x^2-x$, difference $-x^2+2x$.)

\textbf{(b)} Columns 1 and 2 are equal, so $T(x-1)=T(x)-T(1)=0$: $\nul T\supseteq\spn(x-1)$. Columns 1 and 3 are independent, so $\rank T=2$ and $\nul T=\spn(x-1)$. $\range T=\spn\big(1,\ -x^2+2x\big)$. Not an isomorphism.
(Trap: the formula looks invertible. It is not; test injectivity on a basis.)

\textbf{(c)} $1+2x-x^2=T(x^2)$, so the solutions are $x^2+\nul T=\{x^2+c(x-1):c\in\RR\}$.

\subsection*{Question 4}
\textbf{(a)} $v=(v-Pv)+Pv$ with $P(v-Pv)=Pv-P^2v=0$, so $V=\nul P+\range P$. If $v=Pw$ and $Pv=0$, then $v=Pw=P^2w=Pv=0$. So the sum is direct.

\textbf{(b)} ($\Rightarrow$) If $Pv=0$, then $P(Tv)=T(Pv)=0$. If $v=Pw$, then $Tv=TPw=P(Tw)\in\range P$.
($\Leftarrow$) Write $v=n+r$ with $n\in\nul P$, $r\in\range P$; then $Pr=r$. $TPv=Tr$. $PTv=P(Tn)+P(Tr)=0+Tr$, since $Tn\in\nul P$ and $Tr\in\range P$ is fixed by $P$. So $TP=PT$.

\textbf{(c)} $\mathcal C$ contains $0$, and if $S,T\in\mathcal C$ then $(\lambda S+T)P=\lambda SP+TP=\lambda PS+PT=P(\lambda S+T)$.
Take a basis $r_1,\dots,r_k$ of $\range P$ and $n_1,\dots,n_{n-k}$ of $\nul P$; together a basis of $V$ by (a). By (b), $T\in\mathcal C$ exactly when $\MM(T)$ in this basis is block diagonal with a $k\times k$ block and an $(n-k)\times(n-k)$ block. Since $T\mapsto\MM(T)$ is an isomorphism, $\dim\mathcal C=k^2+(n-k)^2$.

\textbf{(d)} $k=2$, $n=3$: $\dim\mathcal C=4+1=5$. Not in $\mathcal C$: $T(x,y,z)=(z,0,0)$ sends $e_3\in\nul P$ to $e_1\notin\nul P$.

\end{document}
